\documentclass[10pt,a4paper]{amsart}
\usepackage[margin=19mm]{geometry}
\usepackage{amsmath,amssymb,mathtools}
\usepackage[scr=rsfs,cal=euler]{mathalfa}
\usepackage{xfrac}
\usepackage[unicode,hidelinks,bookmarks=false]{hyperref}
\newcommand{\ie}{i.e.\ }
\newcommand{\cf}{cf.\ }
\newcommand{\resp}{respectively\ }
\newcommand{\Subsection}{\subsection}
\providecommand{\nobreakdash}{\nobreak\hbox{-}}
\setlength{\emergencystretch}{2em}
\multlinegap0pt
\usepackage{amssymb}
\usepackage[shortlabels]{enumitem}
\usepackage{float}
\newtheorem{lemma}{Lemma}
\title{Exact and Asymptotic Counts of MSTD, MDTS, and Balanced Sets in Dicyclic Groups}
\author{Sagar Mandal, Neetu}
\date{Source: arXiv:2602.09073v1 (CC BY 4.0)}
\begin{document}
\maketitle

\noindent\emph{Source and licence.} Exact and Asymptotic Counts of MSTD, MDTS, and Balanced Sets in Dicyclic Groups. Version arXiv:2602.09073v1 (CC BY 4.0), distributed by arXiv under the Creative Commons Attribution 4.0 International licence, http://creativecommons.org/licenses/by/4.0/, as recorded in the arXiv licence metadata for this exact version and shown on the abstract page https://arxiv.org/abs/2602.09073v1. Full licence notice: \texttt{licenses/arxiv-2602.09073-CC-BY-4.0.txt}.\par\smallskip

\noindent\textit{Editorial scope.} Counted unit: the article's lemma Lemma1.5 (source0.tex lines 534--538). The article's complete original proof of it is delivered with it (source0.tex lines 539--653, one \texttt{proof} environment of 115 lines). No mathematics is written, repaired or re-derived: the statement and the proof are the article's own lines, in the article's own notation, and no retained line is substituted or re-flowed.  No further source-local context is needed: every cross-reference of every retained line resolves inside the two delivered spans.  Nothing was omitted from any retained span, so the package carries no omission record.  No retained line contains a citation, so the wrapper carries no bibliography and no bibliography entry.  No retained line contains a figure environment, an image-inclusion command or drawing code, so no image is delivered and the package declares no asset dependency.  Editorial formatting only, in addition: an AMS \texttt{amsart} preamble that restates the article's own declarations and macros used by the retained lines, namely the environments \texttt{lemma} on separate counters, together with the standard packages those lines need.  The retained environments print in this document as Lemma 1 in order of appearance, whereas the article prints the same items with the numbering of its own full document.  The complete native source of the article as distributed in the arXiv e-print for this exact version is bundled unchanged as \texttt{sources/194293/source0.tex}.
\begin{lemma}\label{Lemma1.5}
Let $n\geq3$ be an odd integer and let 
$$B=\{a^i,a^j,a^kb\}\subset \{1,a,a^{2},\ldots,a^{2n-1},\, b,ab,a^{2}b,\ldots,a^{2n-1}b\},$$ where $0\leq i<j\leq 2n-1,~0\leq k\leq 2n-1$. Then the number of sets of type $B$ that are MSTD, MDTS, and balanced is given as follows.
$$\mathrm{S_{\mathrm{Dic_{4n}}}(3,B)}=2n(n-1)(2n-5),\qquad\mathrm{D_{\mathrm{Dic_{4n}}}(3,B)}=2n(n-1),\qquad\mathrm{B_{\mathrm{Dic_{4n}}}(3,B)}=2n(5n-4).$$
\end{lemma}

\begin{proof}
Let $B=\{a^i,a^j,a^kb\}$ with $0\leq i<j\leq 2n-1 \text{ and }~0\leq k\leq 2n-1$. Then 
$$B+B=\{a^{2i},a^{i+j},a^{2j},a^n\}\cup\{a^{i+k}b,a^{j+k}b,a^{k-i}b,a^{k-j}b\}$$  
$$B-B=\{1,a^{i-j},a^{j-i}\}\cup \{a^{n+i+k}b,a^{n+j+k}b,a^{k+i}b,a^{k+j}b\}.$$ 
Thus, in the absence of any collisions, we have $|B+B|>|B-B|$.
It is straightforward to verify that collisions in $B+B \text{ and }B-B$ occur precisely in the following cases:
\begin{align*}
&\text{(1)}\; i \equiv j \pmod{n}, \quad
&\text{(2)}\; i+j \equiv n \pmod{2n}, \quad
&\text{(3)}\; 2i \equiv 0 \pmod{2n},\\
&\text{(4)}\; i+j \equiv 0 \pmod{2n}, \quad
&\text{(5)}\; 2j \equiv 0 \pmod{2n}, \quad
&\text{(6)}\; n+i \equiv j \pmod{2n},\\
&\text{(7)}\; n+j \equiv i \pmod{2n}.
\end{align*}
Cases 6 and 7 are equivalent to Case 1. Note that the collisions do not depend on the value of $0\leq k\leq 2n-1$. Let $B_{r}$ denote the subset type corresponding to Case $r$, for $r\in \{1,2,3,4,5\}$.   We examine each case individually.\par
\textbf{Case 1}($i \equiv j \pmod n$).
In this case, we necessarily have $(i,j)\in \{(i,i+n): 0\leq i\leq n-1\}$. Consider subsets of the form
$$B_1=\{a^{i},a^{i+n},a^{k}b\}.$$

We first determine when $B_1$ is balanced. This occurs only in the case $(i,j)=(0,n)$. Indeed, when $i=0$, we obtain
$$B_1+B_1=\{1,a^{n}\}\cup\{a^{k}b,a^{n+k}b\},$$
$$B_1-B_1=\{1,a^{n}\}\cup\{a^{n+k}b,a^{k}b\}.$$
Thus $|B_1+B_1|=|B_1-B_1|$, and hence $B_1$ is a balanced set. Since $k$ ranges over $0\le k\le 2n-1$, this yields exactly $2n$ balanced sets.

Now assume $(i,j)\in \{(i,i+n): 1\leq i\leq n-1\}$. In this case,
$$B_1+B_1=\{a^{2i},a^{2i+n},a^{n}\}\cup\{a^{i+k}b,a^{i+n+k}b,a^{k-i}b,a^{k-i-n}b\},$$
whereas
$$B_1-B_1=\{1,a^{n}\}\cup\{a^{n+i+k}b,a^{i+k}b\}.
$$
Consequently, $|B_1+B_1|>|B_1-B_1|$, and thus $B_1$ is an MSTD set. Therefore, no MDTS sets arise in this case.

Finally, for each $i$ with $1\le i\le n-1$ and each $k$ with $0\le k\le 2n-1$, the corresponding set $B_1$ is MSTD. Hence the total number of MSTD sets in this case is
$$\sum_{i=1}^{n-1}2n=2n(n-1).$$
Therefore, in this case, we have 
\begin{equation}\label{eq:case1-counts}
\mathrm{S}_{\mathrm{Dic}_{4n}}(3,B_1)=2n(n-1), \qquad
\mathrm{D}_{\mathrm{Dic}_{4n}}(3,B_1)=0, \qquad
\mathrm{B}_{\mathrm{Dic}_{4n}}(3,B_1)=2n.
\end{equation}
\textbf{Case 2 ($i+j\equiv n \pmod{2n}$).} This implies 
$$(i,j)\in\{(i,n-i): 0\le i\le \frac{n-1}{2}\},$$
and $$(i,j)\in \{(i,3n-i) :n+1\le i\le \frac{3n-1}{2}\}.$$
Since the case $(i,j)=(0,n)$ has already been treated in Case~1, we exclude it from further consideration. For $(i,j)\in\{(i,n-i): 1\le i\le \frac{n-1}{2}\}\cup \{(i,3n-i) :n+1\le i\le \frac{3n-1}{2}\}$, we get 
$$B_2+B_2=\{a^{2i},a^{n},a^{2n-2i}\}\cup
\{a^{i+k}b,a^{n-i+k}b,a^{k-i}b,a^{k-n+i}b\},
$$
and 
$$B_2-B_2=\{1,a^{2i-n},a^{n-2i}\}\cup
\{a^{n+i+k}b,a^{k-i}b,a^{k+i}b,a^{k+n-i}b\}.
$$ In this case we have $|B_2+B_2|=|B_2-B_2|$, and hence every such set $B_2$ is balanced.  Therefore, the total number of balanced sets for each $0\leq k \leq 2n-1$ is
$$\sum_{i=1}^{\frac{n-1}{2}}2n+\sum_{i=n+1}^{\frac{3n-1}{2}}2n
=2n(n-1).
$$
Thus, in this case, we get 
\begin{equation}\label{eq:case2-counts}
\mathrm{S}_{\mathrm{Dic}_{4n}}(3,B_2)=0, \qquad
\mathrm{D}_{\mathrm{Dic}_{4n}}(3,B_2)=0, \qquad
\mathrm{B}_{\mathrm{Dic}_{4n}}(3,B_2)=2n(n-1).
\end{equation}
\textbf{Case 3}$(2i\equiv 0 \pmod{2n})$. 
In this case, $i\in \{0,n\}$. Hence  $$(i,j)\in \{(0,j):1\leq j\leq 2n-1\},$$ or  $$(i,j)\in\{(n,j): n+1\leq j\leq 2n-1\}.$$ Since the case $(i,j)=(0,n)$ has already been analyzed in Case~1,we exclude it from further consideration.  
If $(i,j)\in \{(0,j):1\leq j\leq 2n-1,j\neq n\},$ then  
$$B_{3}+B_{3}=\{1,a^{j},a^{2j},a^n\}\cup\{a^{k}b,a^{j+k}b,a^{k-j}b\},$$ 
$$B_{3}-B_{3}=\{1,a^{2n-j},a^{j}\}\cup \{a^{n+k}b,a^{n+j+k}b,a^{k}b,a^{k+j}b\}.$$
If $(i,j)\in\{(n,j): n+1\leq j\leq 2n-1\}$, then
$$B_{3}+B_{3}=\{1,a^{n+j},a^{2j},a^n\}\cup\{a^{n+k}b,a^{j+k}b,a^{k-j}b\}$$ 
$$B_{3}-B_{3}=\{1,a^{n-j},a^{j-n}\}\cup \{a^{k}b,a^{n+j+k}b,a^{k+n}b,a^{k+j}b\}.$$
In both cases, $|B_3+B_3|=|B_3-B_3|$, and hence every such set $B_3$ is balanced. Therefore, no MSTD or MDTS sets arise in this case. 
For each admissible value of $j$ and for each $k$ with $0\le k\le 2n-1$, we obtain a balanced set. Thus, the total number of balanced sets is
$$
\sum_{\substack{j=1 \\ j\neq n}}^{2n-1}2n
+\sum_{j=n+1}^{2n-1}2n
=6n(n-1).
$$
Therefore, in this case, we obtain
\begin{equation}\label{eq:case3-counts}
\mathrm{S}_{\mathrm{Dic}_{4n}}(3,B_3)=0, \qquad 
\mathrm{D}_{\mathrm{Dic}_{4n}}(3,B_3)=0, \qquad 
\mathrm{B}_{\mathrm{Dic}_{4n}}(3,B_3)=6n(n-1).
\end{equation}\par
\textbf{Case 4}$(i+j\equiv 0 \pmod{2n})$. In this case, we have  $(i,j)\in \{(i,2n-i): 1\leq i\leq n-1\}$. For such pairs $(i,j)$,  we compute 
$$B_{4}+B_{4}=\{1,a^{2i},a^{2n-2i},a^n\}\cup\{a^{i+k}b,a^{k-i}b\},$$ 
$$B_{4}-B_{4}=\{1,a^{2i},a^{2n-2i}\}\cup \{a^{n+i+k}b,a^{n-i+k}b,a^{k+i}b,a^{k-i}b\}.$$
Here, we have $|B_{4}-B_{4}|>|B_{4}+B_{4}|$, and hence every such set $B_4$ is an MDTS set. Consequently, no MSTD or balanced sets arise in this case. For each $i$ with $1\le i\le n-1$ and each $k$ with $0\le k\le 2n-1$, the total number of MDTS sets is
$$
\sum_{i=1}^{n-1}2n=2n(n-1).
$$ Therefore, we obtain \begin{equation}\label{eq:case4-counts}
\mathrm{S}_{\mathrm{Dic}_{4n}}(3,B_4)=0, \qquad 
\mathrm{D}_{\mathrm{Dic}_{4n}}(3,B_4)=2n(n-1), \qquad 
\mathrm{B}_{\mathrm{Dic}_{4n}}(3,B_4)=0.
\end{equation}\par
\textbf{Case 5}($2j\equiv 0 \pmod{2n}$). In this case, we have  $(i,j)\in\{(i,n):0\leq i \leq n-1\}$. Since the case $(i,j)=(0,n)$ has already been analyzed,  we exclude it from further consideration. For pairs $(i,j)\in\{(i,n):1\leq i \leq n-1\}$, we compute
$$B_5+B_5=\{a^{2i},a^{i+n},1,a^{n}\}\cup
\{a^{i+k}b,a^{n+k}b,a^{k-i}b\},$$
$$B_5-B_5=\{1,a^{i-n},a^{n-i}\}\cup
\{a^{n+i+k}b,a^{k}b,a^{k+i}b,a^{k+n}b\}.$$
Here, we have $|B_5+B_5|=|B_5-B_5|$, and hence every such set $B_5$ is balanced. Consequently, no MSTD or MDTS sets arise in this case. For each $i$ with $1\le i\le n-1$ and for each $k$ with $0\le k\le 2n-1$,  total number of balanced sets is
$$\sum_{i=1}^{n-1}2n=2n(n-1).$$
Therefore, in this case we obtain
\begin{equation}\label{eq:case5-counts}
\mathrm{S}_{\mathrm{Dic}_{4n}}(3,B_5)=0, \qquad 
\mathrm{D}_{\mathrm{Dic}_{4n}}(3,B_5)=0, \qquad 
\mathrm{B}_{\mathrm{Dic}_{4n}}(3,B_5)=2n(n-1).
\end{equation}
Therefore, combining \eqref{eq:case1-counts}, \eqref{eq:case2-counts}, \eqref{eq:case3-counts}, \eqref{eq:case4-counts} and \eqref{eq:case5-counts} we get 
\begin{align*}
\mathrm{D_{\mathrm{Dic_{4n}}}(3,B)}&=\mathrm{D_{\mathrm{Dic_{4n}}}(3,B_1)}+\mathrm{D_{\mathrm{Dic_{4n}}}(3,B_2)}+\mathrm{D_{\mathrm{Dic_{4n}}}(3,B_3)}+\mathrm{D_{\mathrm{Dic_{4n}}}(3,B_4)}+\mathrm{D_{\mathrm{Dic_{4n}}}(3,B_5)}\\&=0+0+0+2n(n-1)+0=2n(n-1),\end{align*}
\begin{align*}
\mathrm{B_{\mathrm{Dic_{4n}}}(3,B)}&=\mathrm{B_{\mathrm{Dic_{4n}}}(3,B_1)}+\mathrm{B_{\mathrm{Dic_{4n}}}(3,B_2)}+\mathrm{B_{\mathrm{Dic_{4n}}}(3,B_3)}+\mathrm{B_{\mathrm{Dic_{4n}}}(3,B_4)}+\mathrm{B_{\mathrm{Dic_{4n}}}(3,B_5)}\\&=2n+2n(n-1)+6n(n-1)+0+2n(n-1)=2n(5n-4).\end{align*}
Since the total number of subsets of type $B$ is $\binom{2n}{2}\times 2n$. Hence, 
\begin{align*}
    \mathrm{S}_{\mathrm{Dic}_{4n}}(3,B)&= \binom{2n}{2}\times 2n- \mathrm{D_{\mathrm{Dic_{4n}}}(3,B)}-\mathrm{B_{\mathrm{Dic_{4n}}}(3,B)}\\&=2n^{2}(2n-1)-2n(n-1)-2n(5n-4)=2n(n-1)(2n-5).
\end{align*}
\end{proof}

\end{document}
